
\begin{tikzpicture}[scale=1.0,>=Stealth,font=\scriptsize,line join=round]
%========== (a) sliding screw: flank on flank ==========
\begin{scope}[scale=1.25,yshift=0.12cm]
\fill[gray!30] (0,0) -- (2.4,0) -- (2.4,0.42) -- (0,0.42) -- cycle;
\draw[thick] (0,0) rectangle (2.4,0.42);
\fill[gray!30] (0.45,0.42) -- (0.75,1.30) -- (1.55,1.30) -- (1.85,0.42) -- cycle;
\draw[thick] (0.45,0.42) -- (0.75,1.30) -- (1.55,1.30) -- (1.85,0.42);
\fill[gray!52] (1.62,0.46) -- (1.92,1.34) -- (2.62,1.34) -- (2.62,1.70) -- (0.62,1.70)
  -- (0.62,1.34) -- (0.92,1.34) -- (1.22,0.46) -- cycle;
\draw[thick] (1.62,0.46) -- (1.92,1.34) -- (2.62,1.34) -- (2.62,1.70) -- (0.62,1.70)
  -- (0.62,1.34) -- (0.92,1.34) -- (1.22,0.46);
\draw[red!75!black,line width=1.4pt] (1.58,0.47) -- (1.87,1.32);
\draw[<->,red!75!black,thick] (2.00,0.70) -- (2.28,1.52);
\node[red!75!black,anchor=west,align=left,font=\tiny,text width=2.3cm] at (2.34,1.05)
  {flank \emph{slides} on flank\\ $f \approx 0.15$};
\node[anchor=north,font=\bfseries] at (1.3,-0.56) {(a) sliding screw};
\node[anchor=north,align=center,font=\tiny,text width=3.4cm] at (1.3,-0.98)
  {all the contact is rubbing --- and that friction is what makes it self-locking};
\end{scope}
%========== (b) ball screw: recirculating rolling contact ==========
\begin{scope}[xshift=7.2cm]
% screw shaft with grooves
\begin{scope}
\clip (0,-0.55) rectangle (9.0,0.45);
\fill[gray!30] (0,-0.55) rectangle (9.0,0.45);
\foreach \x in {0.75,1.6,2.45,3.3,4.15,5.0,5.85,6.7,7.55,8.4}{
  \fill[white] (\x,0.45) circle (0.27); \draw[gray!55!black,line width=0.4pt] (\x,0.45) circle (0.27);}
\end{scope}
\draw[thick] (0,-0.55) rectangle (9.0,0.45);
\node[anchor=north,font=\tiny] at (0.55,-0.62) {screw};
% nut body with matching grooves
\begin{scope}
\clip (2.1,0.45) rectangle (6.3,1.42);
\fill[gray!50] (2.1,0.45) rectangle (6.3,1.42);
\foreach \x in {2.45,3.3,4.15,5.0,5.85}{
  \fill[white] (\x,0.45) circle (0.27); \draw[gray!55!black,line width=0.4pt] (\x,0.45) circle (0.27);}
\end{scope}
\draw[thick] (2.1,0.45) rectangle (6.3,1.42);
\node[anchor=west,font=\tiny] at (6.4,0.72) {nut};
% balls
\foreach \x in {2.45,3.3,4.15,5.0,5.85}{
  \draw[thick,fill=blue!22] (\x,0.45) circle (0.24);
  \draw[blue!55!black,line width=0.35pt] (\x-0.09,0.54) arc (135:195:0.14);}
% rotation of the balls
\foreach \x in {3.3,5.0}{\draw[->,blue!60!black,line width=0.6pt] (\x-0.17,0.62) arc (160:20:0.17);}
% recirculation return tube
\draw[very thick,gray!55!black]
  (5.85,0.45) .. controls (6.95,0.45) and (6.95,2.05) .. (4.20,2.05)
  .. controls (1.55,2.05) and (1.55,0.45) .. (2.45,0.45);
\draw[->,gray!55!black,line width=1.1pt] (4.9,2.05) -- (4.1,2.05);
\node[anchor=south,font=\tiny,gray!40!black] at (4.2,2.10) {recirculation return tube};
\draw[->,blue!60!black] (7.35,1.30) -- (6.05,0.70);
\node[blue!55!black,anchor=west,align=left,font=\tiny,text width=2.0cm] at (7.38,1.32)
  {the balls \emph{roll}\\ $f\approx0.003$};
\node[anchor=north,font=\bfseries] at (4.3,-0.72) {(b) ball screw};
\node[anchor=north,align=center,font=\tiny,text width=6.4cm] at (4.3,-1.18)
  {rolling contact removes almost all the friction --- and with it, the self-locking};
\end{scope}
\end{tikzpicture}
